\section{Quantitative consistency and magic-square extraction}\label{sec:extraction}
In this and the next section Bob's operators are the physical operators on
$H_B$, without transposition. Operators on opposite parties commute. For a
vector $v$, $\Sch(v)$ now denotes its nonincreasing Schmidt sequence, padded
with zeros. Multiset unions are sorted before comparison. All vectors other
than the original strategy vector may be unnormalized. Singular-value
perturbation gives
\begin{equation}\label{eq:singularperturbation}
\|\Sch(v)-\Sch(w)\|_2\leq\|v-w\|.
\end{equation}
Indeed, the matrix singular-value inequality follows by applying the
Hermitian eigenvalue inequality to the self-adjoint block matrices with
off-diagonal blocks $K,K^*$; finite-rank approximation gives the compact
Hilbert--Schmidt case. Sorting nonnegative sequences also cannot increase
their Euclidean distance.

\begin{lemma}[Approximate checked classes]\label{lem:consistency}
For projective measurements, suppose mismatched answer classes at a pair
always lose, with absolute losing mass at most $\eta$ on a vector $v$.
Their class projections satisfy $\|(E_c-F_c)v\|\leq\sqrt\eta$.
For the sign tests of $\FMS$, writing $S_x=E^{T_x}_++E^{T_x}_-$,
$X_x=E^{T_x}_+-E^{T_x}_-$ and similarly $S'_y,Y_y$ for Bob, one has
\[
\|(X_xY_y-\sigma_{xy}S_xS'_y)v\|\leq2\sqrt\eta.
\]
\end{lemma}
\begin{proof}
The squared norm of $(E_c-F_c)v$ is the sum of the probabilities of the
disjoint mismatched-class events involving $c$. For signs, expand in the
orthogonal vectors $E^{T_x}_aF^{U_y}_bv$ with $a,b\in\{+,-\}$.
The coefficient $ab-\sigma_{xy}$ is zero on winning events and has magnitude
two on losing events.
\end{proof}

Left multiplication by a contraction preserves a vector estimate. Right
multiplication needs additional control. The following rule will provide it.
If $s$ are Alice reflections with Bob mirrors $\bar s$ satisfying
$\|(s-\bar s)v\|\leq e$, then a word $w$ of length $l$ mirrors to the
reversed Bob word with error at most $le$. Consequently, for an Alice
operator $R$ with $\|R\|\leq2$ and $\|Rv\|\leq re$,
\begin{equation}\label{eq:context}
\|Rwv\|\leq(r+2l)e.
\end{equation}
To see this, replace the rightmost letter by its mirror, commute that
mirror past the remaining Alice letters, and continue. The resulting Bob
word commutes with $R$ and is unitary. The rule holds with the parties
interchanged; the mirroring assertion also holds for contractions.

\begin{theorem}[Square-root word extraction]\label{thm:wordextraction}
Let $\|\phi\|\leq1$ and let local PVMs have valid magic-square row and column
answers, with absolute losing mass at most $\eta$ on each of the nine pairs.
There are local isometries $V=V_A\otimes V_B$ and an auxiliary vector $j$
such that, for every word $w$ of length at most six in Bob's cell reflections
and cell projections, including the empty word,
\begin{equation}\label{eq:wordextraction}
\bigl\|Vw\phi-(I\otimes\widehat w)(\Phi_4\otimes j)\bigr\|
\leq C_{\rm MS}\sqrt\eta,\qquad C_{\rm MS}=2^{24}.
\end{equation}
Here $\Phi_4=\frac12\sum_{z=0}^3|zz\rangle$, and $\widehat w$ uses the
ideal cells in~\eqref{eq:square}, in the same order. The vector $j$ need not
be normalized. The constant is independent of dimension, also on separable
Hilbert spaces.
\end{theorem}

This statement is proved below for the actual row/column game. For comparison,
the analytic estimate of Wu et al.~\cite[Appendix C, (124)]{WBMS} has a
fourth-root rate in a correlation-error parameter; per-pair loss $\eta$
gives checked correlation errors at most $2\eta$. Coladangelo--Stark's
magic-square statement is~\cite[Theorem 6.9]{CSrobust}, applying their general
Theorem 4.16 to an equation/variable game. Their Definition 4.14 uses trace
norm for the state and squared state-dependent operator distance. We do not
use its asserted linear trace-norm rate: Lemma A.3 there fails already for
$v_t=\sqrt{1-t}|0\rangle+\sqrt t|1\rangle$ and $P=|0\rangle\langle0|$,
since the squared projector defect is $t$ but
$\||v_t\rangle\langle v_t|-P\|_1=2\sqrt t>6t$ for $0<t<1/9$.

\begin{proof}[Proof of Theorem~\ref{thm:wordextraction}]
Let $A_{ij},B_{ij}$ be the cell reflections and put $e=2\sqrt\eta$.
Consistency gives $\|(A_{ij}-B_{ij})\phi\|\leq e$. Alice's rows commute
and have product $I$; Bob's columns commute and have products $I,I,-I$.
Write $T\simeq_r U$ for $\|(T-U)\phi\|\leq re$. Put
\[
a=A_{11},\quad b=A_{12},\quad c=A_{21},\quad d=A_{22},\quad
g=A_{31},\quad h=A_{32},\quad i=A_{33}.
\]
Mirroring two letters to a Bob column, using its parity and mirroring back,
gives
\begin{gather*}
[a,c]\simeq_4 0,\qquad [b,d]\simeq_4 0,\\
g\simeq_3 ac,\qquad h\simeq_3 bd,\qquad i\simeq_3-abcd.
\end{gather*}
Also $[a,b]=[c,d]=[g,h]=0$ and $gh=i$ exactly. Substitution of $h$ on
the right, then of $g$ through its two-letter suffix using~\eqref{eq:context},
gives $gh\simeq_{10}acbd$. Thus $acbd\simeq_{13}-abcd$.
Remove the left $a$ and the right $d$, the latter by~\eqref{eq:context},
to obtain $\{b,c\}\simeq_{15}0$.

Similarly $hg\simeq_{10}bdac$, so $acbd\simeq_{20}bdac$.
Commuting $b,d$ through the suffix $ac$ gives $bdac\simeq_8 dabc$;
exact $[c,d]=0$ followed by commuting $b,d$ through $c$ gives
$abcd\simeq_6 adbc$. Combining these with the $13e$ estimate yields
$(ad+da)bc\simeq_{47}0$. Removing $bc$ costs at most $4e$, so
$\{a,d\}\simeq_{51}0$. Transferring two-letter relations to Bob costs
at most $4e$ more. Therefore $55e$ bounds all the corresponding
commutation and anticommutation defects on both sides.

Consider the order-32 real Pauli group
\[
\begin{split}
G=\langle J,a,b,c,d:\;&J^2=a^2=b^2=c^2=d^2=1,\ J\text{ central},\\
&ad=Jda,\ bc=Jcb,\text{ all other pairs commute}\rangle.
\end{split}
\]
Every element has a unique form $J^za^ub^vc^wd^x$ with binary exponents.
Let $f_A,f_B$ evaluate these forms in the physical generators with $J=-I$.
Sorting a product involves at most eight letters and at most 32 adjacent
exchanges. Each exchange costs at most $(55+16)e$ by~\eqref{eq:context},
and cancellations are exact. In particular, on either party,
\begin{equation}\label{eq:almostrepresentation}
\|(f(g)f(h)-f(gh))\phi\|\leq4096e\qquad(g,h\in G).
\end{equation}
Define isometries
\[
W_Av=32^{-1/2}\sum_{g\in G}|g\rangle\otimes f_A(g)v,
\qquad W_Bv=32^{-1/2}\sum_{g\in G}|g\rangle\otimes f_B(g)v.
\]
On the register use the right regular representation
$(R_hF)(g)=F(gh)$. Squaring and averaging~\eqref{eq:almostrepresentation}
over $g$ gives
\begin{equation}\label{eq:regularintertwining}
\|W_Af_A(h)\phi-R_hW_A\phi\|\leq H,\qquad H=4096e,
\end{equation}
and likewise for Bob. An opposite-party isometry preserves these errors.

Since $f(Jg)=-f(g)$ exactly, each isometry has range in the register's
$J=-1$ subspace. Every representation on this subspace is a multiple of
the real four-dimensional representation $a=XI,b=IX,c=IZ,d=ZI$.
For completeness, diagonalize the commuting $c,d$. The involution $a$
exchanges the two $d$ eigenspaces and preserves $c$; $b$ exchanges the
two $c$ eigenspaces and preserves $d$. The four joint eigenspaces have
equal multiplicity, and their identifications by $a,b$ give precisely
these Pauli actions. This also works for infinite multiplicity.
Compose $W_A,W_B$ with these identifications and call the result $V$.

Put $v=V\phi$. For each of the four generators, consistency and
\eqref{eq:regularintertwining} give
$\|(M_s^A-M_s^B)v\|\leq2H+e$.
The four involutions $M_s^AM_s^B$ commute, with common $+1$ projection
$P_0=|\Phi_4\rangle\langle\Phi_4|\otimes I$. Hence
\[
\|(I-P_0)v\|^2
\leq\frac14\sum_s\|(I-M_s^AM_s^B)v\|^2\leq(2H+e)^2.
\]
Taking $\Phi_4\otimes j=P_0v$ proves the empty-word bound
$\|v-\Phi_4\otimes j\|\leq8193e$.

Every other Alice cell equals, or is within $3e$ of, its normal-form word
of degree at most four. Its extraction error is thus at most
$H+3e+8193e$. Bob's cells follow from consistency and
$(M\otimes I)\Phi_4=(I\otimes M^T)\Phi_4$, at a further cost of at most
$e$. The common bound $2^{14}e$ covers a single reflection or projection
on either side; a cell projection is $(I\pm M)/2$.

For a Bob word, mirror its last letter to Alice at cost at most $e$,
commute it past the Bob prefix, apply the prefix estimate, and use the
Alice single-letter estimate on $\Phi_4\otimes j$. Passing from the
physical isometric action to that target costs at most another $8193e$.
Thus each added letter costs at most $e+8193e+2^{14}e$.
For six letters, including the empty-word error, this is less than
$2^{18}e$, which is covered by $2^{24}\sqrt\eta$.
All estimates used absolute losing masses and $\|\phi\|\leq1$ throughout;
there is no normalization factor or dimension dependence.
\end{proof}

\section{The robust flagged sector law}\label{sec:robustsector}
We write $v\approx w$ in this section when
$\|v-w\|\leq C\sqrt\eta$, for a fixed constant independent of the
strategy and its dimension. The constant may increase at successive
steps: only finitely many contractions, bounded sums and words of bounded
length occur.

\begin{lemma}[Two-way spectral domination]\label{lem:domination}
Suppose local contractions satisfy
$\|v-(A\otimes B)w\|\leq d_1$ and
$\|w-(A'\otimes B')v\|\leq d_2$. Then
\[
\|\Sch(v)-\Sch(w)\|_2\leq\sqrt{d_1^2+d_2^2}.
\]
\end{lemma}
\begin{proof}
Let $u=\Sch(v)$, $t=\Sch(w)$, and let $s,s'$ be the Schmidt sequences
after the two contractions, respectively. Singular values decrease
entrywise under contractions, so $s\leq t$, $s'\leq u$.
Equation~\eqref{eq:singularperturbation} gives
$\|(u-t)_+\|_2\leq d_1$ and $\|(t-u)_+\|_2\leq d_2$.
These positive parts have disjoint supports. The same argument for compact
Hilbert--Schmidt maps follows by finite-rank approximation.
\end{proof}

\begin{theorem}[Robust flagged sector law]\label{thm:robustsector}
There is an absolute constant $C_{\rm F}$ such that every projective
$\FMS$ strategy on a vector $\psi$ of norm at most one, with absolute
losing mass at most $\eta$ on each of its twenty-five pairs, satisfies
\begin{equation}\label{eq:robustsector}
\bigl\|\Sch(m)-\mult{4}{(\sqrt2\Sch(f))}\bigr\|_2
\leq C_{\rm F}\sqrt\eta,
\quad f=E_f^{R_1}\psi,\quad m=\psi-f.
\end{equation}
This holds in finite dimension and on separable spatial Hilbert spaces.
\end{theorem}
\begin{proof}
Put $P=E_f^{R_1}$, $Q=F_f^{C_1}$, and write $P^i,Q^j$ for the other
row and column flags. Let $H_y=E_{r_y}^{R_1}$, $G_x=F_{c_x}^{C_1}$.
Cell projections, denoted $E_{i;j,s},F_{j;i,s}$, sum the nonflag answers
with cell sign $s$. Lemma~\ref{lem:consistency} gives
\begin{gather}\label{eq:flagestimates}
P^i\psi\approx Q^j\psi\approx P\psi,\qquad
E_{i;j,s}\psi\approx F_{j;i,s}\psi,\nonumber\\
Qm\approx P^im\approx Q^jm\approx0,\qquad Pm=0.
\end{gather}
The calibrations and signs give
\begin{equation}\label{eq:calibrationestimates}
S_x\psi\approx(Q+G_x)\psi,\quad
S'_y\psi\approx(P+H_y)\psi,\quad
X_xY_y\psi\approx\sigma_{xy}S_xS'_y\psi.
\end{equation}

\paragraph{Mirrors and the nonflag game.}
If $L=F_{j;i,s}$ and $M=E_{i;j,s}$, then $L\psi\approx M\psi$.
Moreover
\[
PM\psi\approx PL\psi=LP\psi\approx LQ^j\psi=0,
\qquad MP\psi\approx MP^i\psi=0.
\]
Thus $Lm\approx Mm$. Replacing rightmost letters and commuting the
opposite-party mirrors proves this assertion for every bounded-length
Bob word, in reversed Alice order. The corresponding reflection estimate
follows by bounded linear combinations.

Merge each row or column flag into a fixed valid nonflag answer, with
the row-one target outside $\{r_0,r_1\}$ and the column-one target outside
$\{c_0,c_1\}$. The merged measurements are PVMs with valid parities.
For a nonflag losing event $\Lambda$, its amplitude on $m$ differs from
its amplitude on $\psi$ by $\Lambda P\psi\approx\Lambda P^i\psi=0$.
Events involving merged flags have amplitudes $O(\sqrt\eta)$
by~\eqref{eq:flagestimates}. There are a fixed number of outcomes, so
the losing \emph{mass} on $m$ is $O(\eta)$ on every pair.
Theorem~\ref{thm:wordextraction} therefore applies at square-root rate.
An unmerged word differs from the merged word by a bounded sum of terms
containing a column flag. Mirror the suffix to Alice, commute it past
that flag and use $Q^jm\approx0$. Each term is $O(\sqrt\eta)$.
Consequently the extraction applies to the unmerged words below as well.

\paragraph{The needed ideal words.}
Let
\[
\widetilde B_y=F_{3;1,(r_y)_3}F_{2;1,(r_y)_2}F_{1;1,(r_y)_1}.
\]
It mirrors $H_y$ on $m$. Define the physical reflections
\[
A=E_{1;2,+}-E_{1;2,-}+P,\quad
B_{12}=F_{2;1,+}-F_{2;1,-}+Q^2,\quad
B=F_{1;2,+}-F_{1;2,-}+Q.
\]
In the ideal square $A=IX$, $B=IZ$. The exact ideal identities are
\begin{align}\label{eq:idealwords}
\widetilde B_0G_1&=\widetilde B_0G_0B_{12},&
\widetilde B_1G_1&=-\widetilde B_1G_0B_{12},\nonumber\\
\widetilde B_1G_0&=B\widetilde B_0G_0.&&
\end{align}
Indeed, $IXc_0=c_1$, $IXr_y=(-1)^yr_y$, $IZr_0=r_1$ and $IZc_0=c_0$.
The reversed order under Bob's transposed action on $\Phi_4$ is essential.
A column answer projector uses two cell projections because of parity,
so the longest word displayed has length six.
Theorem~\ref{thm:wordextraction} transports these identities to $m$.
For its auxiliary vector $j$ the same theorem and singular-value
perturbation give
\begin{equation}\label{eq:robustnormalizations}
\Sch(m)\approx\mult{4}{\bigl(\Sch(j)/2\bigr)},\qquad
\Sch(\widetilde B_0G_0m)\approx\Sch(j)/(2\sqrt2).
\end{equation}
For the second equality, the product of the rank-one row and column
projections has singular value $|\langle c_0|r_0\rangle|=1/\sqrt2$;
the factor $1/2$ comes from $\Phi_4$.

\paragraph{Frustration on the sign supports.}
Put $z_{xy}=S_xS'_ym$, $u=S_0S'_0\psi$ and
$\Delta_A=2P-I$, $\Delta_B=2Q-I$. Calibration implies
$S_xf\approx S'_yf\approx f$. For example,
$S_xP\psi\approx S_xQ\psi=QS_x\psi\approx Q(Q+G_x)\psi\approx f$.
Using mirrors on $m$ then gives
\begin{equation}\label{eq:zidentification}
z_{xy}\approx\widetilde B_yG_xm,\qquad
Pz_{xy}\approx Qz_{xy}\approx0,\qquad
S_xS'_y\psi\approx f+z_{xy}.
\end{equation}
For $P$ annihilation commute $P$ through the Bob word and use $Pm=0$.
For $Q$ annihilation first mirror the \emph{whole} Bob word to Alice,
then commute $Q$ through that word and use $Qm\approx0$.
The three identities~\eqref{eq:idealwords}, consistency $Am\approx B_{12}m$,
and $Af=f$, $Bf\approx f$ yield
\begin{equation}\label{eq:signidentifications}
S_1S'_0\psi\approx Au,\quad S_0S'_1\psi\approx Bu,\quad
S_1S'_1\psi\approx AB\Delta_Au.
\end{equation}
Here $A$ commutes with $P$ exactly, and
$\Delta_Au\approx\Delta_Bu$ by both annihilations
in~\eqref{eq:zidentification}.

Set $X'=AX_1A$, $Y'=BY_1B$. Substituting~\eqref{eq:signidentifications}
in the four sign tests and conjugating by the indicated reflections gives
\begin{gather}\label{eq:fourfrustration}
X_0Y_0u\approx u,\quad X'Y_0u\approx u,\quad X_0Y'u\approx u,
\quad\Delta_AX'\Delta_AY'u\approx-u.
\end{gather}
Subtract the first two equations and apply $Y_0$, using $Y_0^2u=u$,
to get $(X_0-X')u\approx0$. This also holds on $\Delta_Au$:
replace $\Delta_Au$ by $\Delta_Bu$ and commute $\Delta_B$ past the
Alice operators. Replace $X'$ in the fourth equation and add the third.
To remove $Y'$, note that
\[
Y'^2u=BS'_1Bu\approx BS'_1(S_0S'_1\psi)
=B(S_0S'_1\psi)\approx u.
\]
It follows that
\[
(X_0+\Delta_AX_0\Delta_A)u
=2(PX_0P+P^\perp X_0P^\perp)u\approx0.
\]
Since $u\approx f+z_{00}$ and $Pz_{00}\approx0$, the two orthogonal
range blocks imply $PX_0f\approx0$ and $P^\perp X_0z_{00}\approx0$.
Projecting $X_0Y_0u\approx u$ onto these ranges, and using that $Y_0$
commutes with $P$, proves the two directions
\begin{equation}\label{eq:robustbridge}
X_0Y_0f\approx z_{00},\qquad X_0Y_0z_{00}\approx f.
\end{equation}
For example the $P$ projection identifies $PX_0Y_0z_{00}$ with $f$,
and its complementary component is small by the preceding block estimate;
the other direction follows from the complementary projection.

\paragraph{Spectra.}
The operators $X_0,Y_0$ are contractions, so
Lemma~\ref{lem:domination} gives $\Sch(f)\approx\Sch(z_{00})$.
Combine~\eqref{eq:zidentification} and~\eqref{eq:robustnormalizations}:
four copies of $\Sch(j)/2$ are four copies of
$\sqrt2\,\Sch(j)/(2\sqrt2)$. This proves~\eqref{eq:robustsector}.
All operations have dimension-independent bounds and bounded word lengths.
They establish a finite absolute $C_{\rm F}$; no useful numerical value
for this constant is asserted.
\end{proof}

\section{Gluing estimates and logarithmic rank}\label{sec:rank}
\begin{lemma}[Approximate gluing]\label{lem:robustgluing}
There is a constant $C_{\rm G}$ such that a unit-vector PVM strategy for
$\Gamma$ of loss at most $\eps$ has the following properties.
Let $L_A,L_B$ be Bob's two label-answer projections, put
\[
\alpha=L_A\psi,\quad\beta=L_B\psi,\quad
\phi_2=E_f^{(2,R_1)}\psi,\qquad
c=\Sch(\alpha),\quad b=\Sch(\beta),\quad a=\Sch(\phi_2).
\]
The subscripts on $L_A,L_B$ denote answers, not tensor factors.
With $P_\bot=E_\bot^{(2,R_1)}$, $Q_\bot=F_\bot^{(2,C_1)}$,
$t=\Sch(P_\bot Q_\bot\psi)$ and $\delta=C_{\rm G}\sqrt\eps$,
\begin{align}\label{eq:fivegluing}
\|c-(a\sq t)\|_2&\leq\delta,\tag{i}\\
\|c-\mult{4}{(\sqrt2 b)}\|_2&\leq\delta,\tag{ii}\\
\|b-\mult{4}{(\sqrt2 a)}\|_2&\leq\delta,\tag{iii}\\
\|\Sch(\psi)-(c\sq b)\|_2&\leq\delta,\tag{iv}\\
\|\alpha\|^2+\|\beta\|^2&=1.\tag{v}
\end{align}
\end{lemma}
\begin{proof}
Every supported pair has losing mass at most $52\eps$.
Link one identifies the first-copy magic-square vector with $\alpha$
and its flag vector with $\beta$, with vector error $O(\sqrt\eps)$.
Apply Theorem~\ref{thm:robustsector} and singular-value perturbation
to obtain (ii).

For (iii) remove the boundary coherently: put
$\psi'=(I-P_\bot)\psi$. Boundary consistency on all twenty-five
copy-two pairs identifies each boundary vector with $P_\bot\psi$
within $O(\sqrt\eps)$. In particular every row or column boundary
projection applied to $\psi'$ is $O(\sqrt\eps)$.
Merge the boundary into a valid retained answer, choosing a magic-square
answer at a row or column and an ordinary triple at a sign question.
For a nonboundary losing projector $\Lambda$,
$\Lambda P_\bot\psi$ is small: replace $P_\bot$ by the boundary
projection at the relevant opposite-party question, where its product
with $\Lambda$ is zero. Thus every losing amplitude on $\psi'$ is
$O(\sqrt\eps)$ and the merged FMS losing mass is $O(\eps)$.
At $R_1$ its flag vector is exactly $\phi_2$, and its nonflag vector is
exactly $E_{\rm MS}^{(2,R_1)}\psi$, which link two identifies with
$\beta$. The robust sector law proves (iii).

For (i), link two identifies $\alpha$ with
$(E_f^{(2,R_1)}+P_\bot)\psi$. At the pair $((2,R_1),(2,C_1))$ the
three classes flag, boundary and magic square are consistent, so this
vector is within $O(\sqrt\eps)$ of
\[
w=E_f^{(2,R_1)}F_f^{(2,C_1)}\psi+P_\bot Q_\bot\psi.
\]
Its two summands are orthogonal blocks on \emph{both} parties, hence
their Schmidt lists form the union $\Sch(w)$. The first block is
within $O(\sqrt\eps)$ of $\phi_2$. Singular-value perturbation and
sorting prove (i). For (iv) use the two-sided block vector
\[
E_{\rm MS}^{(1,R_1)}L_A\psi+E_f^{(1,R_1)}L_B\psi.
\]
Consistency makes this vector close to $\psi$ and its blocks close to
$\alpha,\beta$, respectively. Finally (v) is exact completeness of
the label PVM. A single enlarged constant covers all four errors.
\end{proof}

\begin{lemma}[Descent with multiplicities]\label{lem:descent}
Suppose finitely supported nonnegative sorted sequences $a,b,c,t$
satisfy (i)--(iii) in the infinity norm with error at most $\delta>0$.
Let $\mu_0=\max a$. For every integer $k\geq0$ with
$\mu_02^{-k}\geq64\delta$ there are distinct positive entries
$\mu_k=a_{i_k}$ such that
\begin{equation}\label{eq:descenterror}
|\mu_k-\mu_02^{-k}|\leq(2+\sqrt2)\delta.
\end{equation}
If $K$ is the number of these integers, then
$K\geq\log_2(\mu_0/(64\delta))$.
At each index $i_k$, four entries of $b$ are within $\delta$ of
$\sqrt2\mu_k$ and sixteen entries of $c$ are within
$(1+\sqrt2)\delta$ of $2\mu_k$. These are disjoint index blocks
for distinct $k$.
\end{lemma}
\begin{proof}
For a positive entry $\mu$ of $a$, select its occurrence $p$ in the
sorted union $a\sq t$, retaining occurrences when values coincide.
Read (i) and (ii) at $p$ to obtain
$|b_{\lceil p/4\rceil}-\mu/\sqrt2|\leq\sqrt2\delta$.
Read (iii) at $\lceil p/4\rceil$ to find an entry of the zero-padded
$a$ within
$g_0=(1+1/\sqrt2)\delta$ of $\mu/2$.
Iteration contracts the preceding error by $1/2$, so its total is at
most $g_1=2g_0=(2+\sqrt2)\delta$.
Retain exactly the integers in the statement. Their entries are
positive, so none is a padded zero. Consecutive entries differ by at
least $\mu_02^{-k}/2-2g_1>0$. Thus their $a$ indices are distinct.
If $\mu_0\geq64\delta$, the number retained is
$\lfloor\log_2(\mu_0/(64\delta))\rfloor+1$; if not, take $K=0$.
Both conventions give the stated inequality.

Relation (iii) supplies the four $b$ indices
$4(i_k-1)+1,\ldots,4i_k$. Relation (ii) supplies the sixteen $c$
indices $16(i_k-1)+1,\ldots,16i_k$, with the stated accumulated error.
Distinct $i_k$ give disjoint blocks, regardless of equal spectral values.
\end{proof}

\begin{proof}[Proof of Theorem~\ref{thm:sharp}]
First consider a finite-dimensional mixed-state POVM strategy on local
spaces of dimensions $d_A,d_B$. Loss is linear in the state, so some
pure component has loss at most $\eps$ using the same measurements.
Let its Schmidt rank be $r\leq\min(d_A,d_B)$.
Give the finitely many local POVMs common Naimark dilations with fixed
pure local ancillas. The resulting PVM strategy preserves the
correlations and the Schmidt rank $r$. The dilated ambient dimensions
may be larger and are never used as the original dimensions. Every
one-sided projection and local contraction used below has Schmidt rank
at most $r$.

Apply Lemma~\ref{lem:robustgluing}. Write
$x=\|a\|_2$, $y=\|b\|_2$, $z=\|c\|_2$.
Its relations give
\[
z\leq\sqrt8y+\delta,\quad y\leq\sqrt8x+\delta,\quad y^2+z^2=1.
\]
Since $y^2+(\sqrt8y+\delta)^2\leq(3y+\delta)^2$,
\begin{equation}\label{eq:sectorweight}
y\geq(1-\delta)/3,\qquad
x\geq(1-4\delta)/\sqrt{72}.
\end{equation}
Choose $\eps_0>0$ so that $C_{\rm G}\sqrt{\eps_0}\leq1/100$.
Then $x>0.11$. As $\operatorname{rank}(\phi_2)\leq r$,
$\mu_0=\max a\geq0.11/\sqrt r$.

The 2-norm estimates imply infinity-norm estimates, so
Lemma~\ref{lem:descent} applies. All its selected $b$ and $c$ entries
exceed $\delta$: even the smallest chain entry is at least
$(64-2-\sqrt2)\delta$.
Relation (iv) then forces at least $20K$ positive entries in
$\Sch(\psi)$. Indeed, match occurrences in the sorted union
$c\sq b$; an occurrence larger than $\delta$ cannot match a zero
with coordinate error at most $\delta$. Multiplicity cannot be lost
by coincident values. Consequently
\begin{align}\label{eq:rankimplicit}
r&\geq20K\geq20\log_2\frac{\mu_0}{64C_{\rm G}\sqrt\eps}\nonumber\\
&\geq10\log_2(1/\eps)-10\log_2r
 -20\log_2(64C_{\rm G}/0.11).
\end{align}
Put $L=\log_2(1/\eps)$. If $r\geq10L$ the desired lower bound is
immediate. Otherwise $\log_2r\leq\log_2L+\log_210$ in
\eqref{eq:rankimplicit}. Absorb the fixed constants to get
$r\geq10L-10\log_2L-C$. This bounds each original local dimension.
Together with Theorem~\ref{thm:main} it proves the assertion for
$D_\Gamma$ and all the stated quantifiers.
At zero loss the same descent has arbitrarily many positive entries;
one treats that case directly, without taking a logarithm of zero.
\end{proof}
